\chapter{16550A 串口：第一个观测通道}

\section{初始化顺序}

固件先关闭 UART 中断，设置 DLAB，写入波特率除数 1，再关闭 DLAB 并配置
8 数据位、无校验、1 停止位。随后配置 FIFO 和 modem control。对 QEMU
模拟串口而言，宿主终端不真正按 115200 波特限制吞吐，但寄存器协议仍应按照
真实设备编写。

\begin{longtable}{p{0.24\textwidth}p{0.20\textwidth}p{0.46\textwidth}}
\toprule
动作 & 端口/值 & 含义 \\
\midrule
关闭中断 & 0x3F9 / 0x00 & 最早阶段采用轮询，不接收 IRQ \\
打开 DLAB & 0x3FB / 0x80 & 让偏移 0、1 指向除数锁存器 \\
设置除数 & 0x3F8 / 1，0x3F9 / 0 & 115200 基准波特率 \\
设置 8N1 & 0x3FB / 0x03 & 同时关闭 DLAB \\
设置 FIFO & 0x3FA / 0xC7 & 启用并清空 FIFO \\
设置 modem & 0x3FC / 0x0B & 拉起 DTR、RTS、OUT2 \\
\bottomrule
\end{longtable}

\section{发送不是直接写字节}

写数据前读取线路状态寄存器，等待 bit 5 表示发送保持寄存器为空。若直接连续
写而不检查，代码隐含了设备永远足够快的假设。当前实现用 16 位计数器进行最多
\code{0xFFFF} 次轮询，成功通过进位标志返回，超时清除进位标志。

\section{两条日志的协议意义}

\code{RESET} 证明复位跳转、栈和 UART 初始化已经工作；
\code{SERIAL_READY} 证明完整字符串发送路径仍然可用。失败镜像在第一条消息
之后人为屏蔽就绪状态，用同一轮询函数走超时路径。因此测试既能识别“完全未
执行”，也能识别“执行了但设备没有完成”。
